\documentclass[a4paper,12pt]{report}

\usepackage[utf8]{inputenc}
\usepackage[T1]{fontenc}
\usepackage[french]{babel}
\usepackage{tikz}
\usetikzlibrary{shapes,arrows,positioning,fit,calc}
\usepackage{amsmath}
\usepackage{booktabs}
\usepackage{graphicx}
\usepackage{geometry}
\usepackage{xcolor}
\usepackage{colortbl}
\geometry{margin=2.5cm}

\begin{document}

\title{Étude théorique du système tracteur-remorque}
\author{Rapport de stage ingénieur en robotique}
\date{\today}

\maketitle

\tableofcontents

\chapter{Étude théorique du système tracteur-remorque}
\label{chap:etude_theorique}

\section*{Introduction}
\addcontentsline{toc}{section}{Introduction}

Ce chapitre présente l'étude théorique du système tracteur-remorque développé dans le cadre de ce stage. L'objectif est de modéliser la cinématique d'un robot différentiel articulé avec une remorque passive, et de concevoir une architecture de fusion de capteurs et de contrôle sécurisé pour la navigation autonome.

La démarche adoptée consiste à valider les modèles théoriques par un prototypage initial sous simulateur Gazebo, utilisant des primitives géométriques simples (boîtes) pour le tracteur et la remorque. Cette approche permet de tester rapidement les algorithmes de fusion de capteurs, d'estimation de l'angle d'attelage et de prévention du jackknife avant de passer à la modélisation CAO détaillée.

\begin{center}
\begin{tikzpicture}[node distance=1.5cm, auto]
    \node[draw, fill=blue!10, rectangle, rounded corners, minimum width=2.5cm, minimum height=1cm] (model) {Modélisation\\cinématique};
    \node[draw, fill=green!10, rectangle, rounded corners, minimum width=2.5cm, minimum height=1cm, right=of model] (stabilite) {Analyse de\\stabilité};
    \node[draw, fill=yellow!10, rectangle, rounded corners, minimum width=2.5cm, minimum height=1cm, right=of stabilite] (ekf) {Fusion EKF\\multi-capteurs};
    \node[draw, fill=orange!10, rectangle, rounded corners, minimum width=2.5cm, minimum height=1cm, right=of ekf] (controle) {Contrôle\\sécurisé};
    \node[draw, fill=red!10, rectangle, rounded corners, minimum width=2.5cm, minimum height=1cm, right=of controle] (proto) {Prototype\\Gazebo};
    
    \draw[->, thick] (model) -- (stabilite);
    \draw[->, thick] (stabilite) -- (ekf);
    \draw[->, thick] (ekf) -- (controle);
    \draw[->, thick] (controle) -- (proto);
\end{tikzpicture}
\end{center}
\vspace{0.5cm}

\section{Modélisation cinématique tracteur-remorque}
\label{sec:modelisation_cinematique}

\subsection{Repères et paramètres géométriques}
\label{subsec:repres_geometriques}

Le système tracteur-remorque est modélisé comme un robot différentiel tracteur tirant une remorque passive reliée par un attelage pivotant (figure \ref{fig:schema_cinematique}). Les repères suivants sont définis :

\begin{itemize}
    \item $\mathcal{R}_t = (O_t, \vec{x}_t, \vec{y}_t)$ : repère lié au tracteur, d'origine au centre de l'essieu arrière
    \item $\mathcal{R}_r = (O_r, \vec{x}_r, \vec{y}_r)$ : repère lié à la remorque, d'origine au centre de son essieu
    \item $\mathcal{R}_g = (O_g, \vec{x}_g, \vec{y}_g)$ : repère global inertiel
\end{itemize}

\begin{figure}[htbp]
    \centering
    \begin{tikzpicture}[scale=1.5]
        % Repere global
        \draw[->, thick] (0,0) -- (3,0) node[right] {$\vec{x}_g$};
        \draw[->, thick] (0,0) -- (0,2) node[above] {$\vec{y}_g$};
        
        % Tracteur
        \draw[fill=blue!20, thick] (1,0.5) rectangle (2.5,1.0);
        \node at (1.75,0.75) {Tracteur};
        \draw[->, thick, red] (1.75,0.75) -- (2.25,0.75) node[right] {$\vec{x}_t$};
        
        % Remorque
        \draw[fill=green!20, thick] (1.5,-0.8) rectangle (3.0,-0.3);
        \node at (2.25,-0.55) {Remorque};
        \draw[->, thick, red] (2.25,-0.55) -- (2.75,-0.55) node[right] {$\vec{x}_r$};
        
        % Attelage
        \draw[dashed, thick] (1.75,0.5) -- (2.25,-0.3);
        \filldraw (1.75,0.5) circle (0.05);
        \filldraw (2.25,-0.3) circle (0.05);
        
        % Beta angle
        \draw[->, thick, purple] (1.75,0.5) -- (1.95,0.2) node[midway, left] {$\beta$};
        
        % Theta angles
        \node at (0.5,1.3) {$\theta_t$};
        \node at (1.0,-0.1) {$\theta_r$};
        
        % Distance d
        \draw[<->] (1.9,0.5) -- (1.9,-0.3) node[midway, right] {$d$};
    \end{tikzpicture}
    \caption{Schéma cinématique du système tracteur-remorque}
    \label{fig:schema_cinematique}
\end{figure}

Les paramètres géométriques du système sont :

\begin{table}[htbp]
    \centering
    \begin{tabular}{lcc}
        \toprule
        Paramètre & Symbole & Valeur (m) \\
        \midrule
        Distance attelage-essieu remorque & $d$ & 0.50 \\
        Entraxe roues tracteur & $L$ & 0.40 \\
        Longueur tracteur & $l_t$ & 0.60 \\
        Largeur tracteur & $w_t$ & 0.40 \\
        Longueur remorque & $l_r$ & 0.80 \\
        Largeur remorque & $w_r$ & 0.50 \\
        \bottomrule
    \end{tabular}
    \caption{Paramètres géométriques du système}
    \label{tab:parametres_geometriques}
\end{table}

L'état du système est défini par :
\begin{equation}
    \mathbf{x} = [x_t, y_t, \theta_t, \theta_r, \beta]^T
    \label{eq:etat_systeme}
\end{equation}
où $(x_t, y_t)$ est la position du tracteur dans le repère global, $\theta_t$ son orientation, $\theta_r$ l'orientation de la remorque, et $\beta = \theta_t - \theta_r$ l'angle d'attelage.

\subsection{Dérivation de l'équation cinématique de l'attelage}
\label{subsec:derivation_beta}

La cinématique de l'angle d'attelage $\beta$ est dérivée en considérant la contrainte géométrique imposée par la liaison pivot entre le tracteur et la remorque. Soit $H$ le point d'attelage sur le tracteur et $P$ le point d'attelage sur la remorque.

La position du point d'attelage dans le repère tracteur est :
\begin{equation}
    \mathbf{H}_t = \begin{bmatrix} -d_h \\ 0 \end{bmatrix}
    \label{eq:position_attelage_tracteur}
\end{equation}
où $d_h$ est la distance du centre de l'essieu au point d'attelage.

La position du même point dans le repère remorque est :
\begin{equation}
    \mathbf{P}_r = \begin{bmatrix} d \\ 0 \end{bmatrix}
    \label{eq:position_attelage_remorque}
\end{equation}
où $d$ est la distance du point d'attelage à l'essieu de la remorque.

La contrainte de liaison s'écrit dans le repère global :
\begin{equation}
    \mathbf{R}(\theta_t) \mathbf{H}_t + \begin{bmatrix} x_t \\ y_t \end{bmatrix} = \mathbf{R}(\theta_r) \mathbf{P}_r + \begin{bmatrix} x_r \\ y_r \end{bmatrix}
    \label{eq:contrainte_liaison}
\end{equation}
où $\mathbf{R}(\theta)$ est la matrice de rotation d'angle $\theta$.

En dérivant cette contrainte par rapport au temps et en utilisant les relations de cinématique du robot différentiel :
\begin{equation}
    \dot{x}_t = v \cos\theta_t, \quad \dot{y}_t = v \sin\theta_t, \quad \dot{\theta}_t = \omega
    \label{eq:cinematique_tracteur}
\end{equation}

On obtient l'équation cinématique de l'angle d'attelage :
\begin{equation}
    \dot{\beta} = \omega - \frac{v}{d} \sin\beta
    \label{eq:beta_dot}
\end{equation}

Cette équation constitue le fondement du modèle cinématique utilisé pour l'estimation et le contrôle.

\subsection{Analyse de stabilité et angle limite de jackknife}
\label{subsec:angle_limite_jackknife}

Le phénomène de jackknife correspond à une instabilité où la remorque se plie excessivement par rapport au tracteur, pouvant atteindre la limite mécanique de l'attelage. L'angle limite $\beta_{\text{lim}}$ est déterminé par la géométrie du système.

Pour un virage de rayon de courbure $R$, la relation cinématique en régime permanent ($\dot{\beta} = 0$) donne :
\begin{equation}
    \omega = \frac{v}{d} \sin\beta_{\text{ss}}
    \label{eq:regime_permanent}
\end{equation}

En utilisant la relation $R = v/\omega$ pour le tracteur, on obtient l'angle d'attelage en régime permanent :
\begin{equation}
    \beta_{\text{ss}} = \arcsin\left(\frac{d}{R}\right)
    \label{eq:beta_steady_state}
\end{equation}

Le rayon de virage minimal $R_{\min}$ correspondant à l'angle limite mécanique $\beta_{\max}$ est donc :
\begin{equation}
    R_{\min} = \frac{d}{\sin\beta_{\max}}
    \label{eq:rayon_minimal}
\end{equation}

Avec les paramètres du système ($d = 0.50$ m, $\beta_{\max} = 45^\circ$), on obtient :
\begin{equation}
    R_{\min} = \frac{0.50}{\sin(45^\circ)} \approx 0.71~\text{m}
    \label{eq:rayon_minimal_calcule}
\end{equation}

Cette valeur est publiée sur le topic \texttt{/trailer/r\_min} par le nœud EKF.

La condition de stabilité du système peut être analysée en linéarisant l'équation (\ref{eq:beta_dot}) autour de $\beta = 0$ :
\begin{equation}
    \delta\dot{\beta} = \omega - \frac{v}{d} \delta\beta
    \label{eq:linearisation_beta}
\end{equation}

Le pôle du système linéarisé est $-\frac{v}{d}$, ce qui indique une stabilité asymptotique pour $v > 0$ (marche avant) et une instabilité pour $v < 0$ (marche arrière). Cela justifie la nécessité d'une stratégie de contrôle spécifique pour le recul.

\section{Architecture de fusion de capteurs par EKF}
\label{sec:architecture_ekf}

\subsection{Principe de la fusion multi-capteurs}
\label{subsec:principe_fusion}

L'estimation de la pose du robot et de l'angle d'attelage repose sur une architecture de filtre de Kalman étendu (EKF) à trois niveaux, permettant de fusionner des capteurs de fréquences et de précisions différentes.

\begin{table}[htbp]
    \centering
    \begin{tabular}{lccc}
        \toprule
        Capteur & Fréquence & Type de donnée & Rôle dans la fusion \\
        \midrule
        Odométrie roues & 50 Hz & $v, \omega$ & Estimation locale haute fréquence \\
        IMU & 50 Hz & $\dot{\theta}$ & Stabilisation angulaire \\
        GPS & 10 Hz & $x, y$ (LLA) & Correction position globale \\
        Joint hitch & 50 Hz & $\beta_{\text{joint}}$ & Mesure directe attelage \\
        \bottomrule
    \end{tabular}
    \caption{Capteurs utilisés dans l'architecture EKF}
    \label{tab:capteurs_ekf}
\end{table}

\subsection{Architecture à trois niveaux}
\label{subsec:architecture_3_niveaux}

L'architecture EKF (figure \ref{fig:architecture_ekf}) est organisée en trois niveaux :

\begin{enumerate}
    \item \textbf{EKF Local (niveau 1)} : Fusionne l'odométrie des roues et l'IMU pour produire une estimation locale continue de la pose dans le repère \texttt{odom}. Ce filtre ne reçoit jamais de correction GPS directe, conformément à la REP 105.
    
    \item \textbf{NavSat Transform (niveau 2)} : Convertit les coordonnées GPS (Latitude/Longitude/Altitude) en coordonnées cartésiennes planes $(x, y)$ par rapport à un datum fixe. Il utilise l'orientation de l'odométrie pour éviter les ambiguïtés de yaw du GPS.
    
    \item \textbf{EKF Global (niveau 3)} : Fusionne l'odométrie, l'IMU et la position GPS convertie pour produire la pose globale dans le repère \texttt{map}, corrigeant la dérive à long terme de l'odométrie.
\end{enumerate}

\begin{figure}[htbp]
    \centering
    \begin{tikzpicture}[scale=0.9, every node/.style={transform shape}]
        % Capteurs
        \node[draw, fill=blue!10, rectangle, minimum width=2cm, minimum height=1cm] (odom) at (0, 4) {Odométrie};
        \node[draw, fill=blue!10, rectangle, minimum width=2cm, minimum height=1cm] (imu) at (0, 2.5) {IMU};
        \node[draw, fill=blue!10, rectangle, minimum width=2cm, minimum height=1cm] (gps) at (0, 1) {GPS};
        
        % EKF Local
        \node[draw, fill=green!20, rectangle, minimum width=2.5cm, minimum height=1.5cm] (ekf_local) at (4, 3.25) {EKF Local\\(30 Hz)};
        
        % NavSat Transform
        \node[draw, fill=yellow!20, rectangle, minimum width=2.5cm, minimum height=1.5cm] (navsat) at (4, 1) {NavSat\\Transform};
        
        % EKF Global
        \node[draw, fill=red!20, rectangle, minimum width=2.5cm, minimum height=1.5cm] (ekf_global) at (8, 2.25) {EKF Global\\(30 Hz)};
        
        % TF outputs
        \node[draw, fill=purple!20, rectangle, minimum width=2cm, minimum height=1cm] (tf_odom) at (12, 4) {TF: odom$\to$base};
        \node[draw, fill=purple!20, rectangle, minimum width=2cm, minimum height=1cm] (tf_map) at (12, 2.25) {TF: map$\to$odom};
        
        % Arrows
        \draw[->, thick] (odom) -- (ekf_local);
        \draw[->, thick] (imu) -- (ekf_local);
        \draw[->, thick] (imu) -- (ekf_global);
        \draw[->, thick] (gps) -- (navsat);
        \draw[->, thick] (navsat) -- (ekf_global);
        \draw[->, thick] (ekf_local) -- (ekf_global);
        \draw[->, thick] (ekf_local) -- (tf_odom);
        \draw[->, thick] (ekf_global) -- (tf_map);
        
        % Labels
        \node at (2, 3.8) {\footnotesize /odom};
        \node at (2, 2.8) {\footnotesize /imu};
        \node at (2, 1.3) {\footnotesize /gps/fix};
        \node at (6, 3) {\footnotesize /odometry/local};
        \node at (6, 1.5) {\footnotesize /odometry/gps};
        \node at (10, 2.5) {\footnotesize /odometry/global};
    \end{tikzpicture}
    \caption{Architecture EKF à trois niveaux pour la fusion multi-capteurs}
    \label{fig:architecture_ekf}
\end{figure}

La configuration de cette architecture est définie avec les fréquences de publication fixées à 30 Hz pour assurer un temps de réponse suffisant pour le contrôle.

\section{Stratégie de recul sécurisé}
\label{sec:strategie_recul}

\subsection{Problématique du recul avec remorque}
\label{subsec:problematique_recul}

L'analyse de stabilité de la section \ref{subsec:angle_limite_jackknife} a montré que le système est instable en marche arrière ($v < 0$). En effet, l'équation (\ref{eq:beta_dot}) montre que pour $v < 0$, le terme $-\frac{v}{d}\sin\beta$ devient positif lorsque $\beta > 0$, créant une rétroaction positive qui amplifie l'angle d'attelage jusqu'au jackknife.

La stratégie de recul sécurisé implémente donc deux mécanismes :
\begin{enumerate}
    \item Une limitation adaptative de la vitesse angulaire $\omega$ en fonction de l'angle $\beta$ courant
    \item Une réduction de la vitesse linéaire lorsque l'angle d'attelage est élevé
\end{enumerate}

\subsection{Logique de contrôle du recul sécurisé}
\label{subsec:logique_recul}

Le tableau \ref{tab:logique_controle} décrit la logique de contrôle implémentée :

\begin{table}[htbp]
    \centering
    \rowcolors{2}{gray!10}{white}
    \begin{tabular}{p{3cm}p{4cm}p{4cm}}
        \toprule
        \textbf{Étape} & \textbf{Condition} & \textbf{Action} \\
        \midrule
        1. Limitation $v$ & Toujours & $v \in [-v_{\max,\text{rev}}, v_{\max,\text{fwd}}]$ \\
        2. Calcul $\omega_{\max}$ & $\beta_{\text{ratio}} = |\beta|/\beta_{\max}$ & $\omega_{\max} = \omega_{\max,\text{stat}}(1-0.35\beta_{\text{ratio}})$ \\
        3. Limite cinématique & $|v| > 0.05$ m/s & $\omega_{\max} = \min(\omega_{\max}, \frac{\sin\beta_{\max}}{d}|v|)$ \\
        4. Réduction arrière & $v < 0$ et $|\beta| > 0.2$ & $v \leftarrow v \times (1-0.5\beta_{\text{ratio}})$ \\
        5. Application & Final & $(v_{\text{safe}}, \omega_{\text{safe}})$ \\
        \bottomrule
    \end{tabular}
    \caption{Logique de contrôle du recul sécurisé}
    \label{tab:logique_controle}
\end{table}

\section{Stack technologique de simulation}
\label{sec:stack_simulation}

\subsection{Outils et frameworks utilisés}
\label{subsec:outils_simulation}

Le tableau \ref{tab:stack_simulation} répertorie l'ensemble des outils technologiques utilisés pour le prototypage initial sous Gazebo.

\begin{table}[htbp]
    \centering
    \begin{tabular}{llp{6cm}}
        \toprule
        Catégorie & Outil/Bibliothèque & Justification \\
        \midrule
        Middleware ROS & ROS2 Humble & Version LTS stable, support Nav2 natif, architecture modulaire \\
        Simulateur & Gazebo Ignition Fortress & Physique réaliste, intégration ROS2 native, plugins capteurs \\
        Navigation & Nav2 & Stack de navigation autonome standard ROS2, planificateurs adaptatifs \\
        Langage & Python 3.10+ & Développement rapide de nœuds de contrôle et EKF \\
        Langage & C++ (ament\_cmake) & Performance pour les nœuds critiques (robot\_state\_publisher) \\
        Bibliothèque math & numpy & Calculs matriciels pour EKF, opérations vectorielles \\
        Bibliothèque TF & tf2\_ros & Gestion de l'arbre des transformations repères \\
        Fusion capteurs & robot\_localization & Package standard pour EKF multi-capteurs \\
        Contrôle & ros2\_control & Interface standard contrôleurs-actionneurs ROS2 \\
        Visualisation & RViz2 & Visualisation TF, costmaps, trajectoires Nav2 \\
        \bottomrule
    \end{tabular}
    \caption{Stack technologique de simulation}
    \label{tab:stack_simulation}
\end{table}

\subsection{Justification des choix technologiques}
\label{subsec:justification_choix}

Le choix de ROS2 Humble plutôt que ROS1 est motivé par la nécessité d'une architecture temps réel déterministe pour le contrôle de sécurité anti-jackknife. Le modèle de communication DDS (Data Distribution Service) de ROS2 offre une meilleure fiabilité et une latence plus prévisible que le système TCP/UDP de ROS1.

Gazebo Ignition Fortress (aujourd'hui Gazebo Sim) a été choisi pour son support natif de ROS2 via les ponts \texttt{ros\_gz\_bridge}, évitant les problèmes de synchronisation d'horloge rencontrés avec Gazebo Classic. Les plugins de capteurs (LIDAR, IMU, GPS) offrent une modélisation réaliste du bruit, essentielle pour valider les algorithmes de fusion.

L'utilisation de Nav2 avec le planificateur SMAC Lattice permet de générer des trajectoires respectant les contraintes cinématiques du robot différentiel, tandis que le contrôleur Regulated Pure Pursuit assure un suivi lissé avec adaptation de vitesse dans les virages.

\section{Architecture logicielle du prototype}
\label{sec:architecture_logicielle}

\subsection{Vue d'ensemble de l'architecture ROS2}
\label{subsec:vue_ensemble_ros2}

L'architecture logicielle (figure \ref{fig:architecture_ros2}) est organisée en nœuds ROS2 communicant via des topics, services et actions. Les nœuds de sécurité s'intercalent entre la stack Nav2 et le contrôleur de base pour garantir le respect des contraintes cinématiques de la remorque.

\begin{figure}[htbp]
    \centering
    \begin{tikzpicture}[scale=0.7, every node/.style={transform shape}]
        % Nav2 Stack
        \node[draw, fill=blue!10, rectangle, minimum width=3cm, minimum height=1.5cm] (nav2) at (0, 5) {Nav2 Stack\\(planner, controller)};
        
        % Trailer-Aware Controller
        \node[draw, fill=green!20, rectangle, minimum width=3cm, minimum height=1.2cm] (trailer_ctrl) at (0, 3) {Trailer-Aware\\Controller\\(10 Hz)};
        
        % CMD Vel Safety
        \node[draw, fill=yellow!20, rectangle, minimum width=3cm, minimum height=1.2cm] (safety) at (0, 1) {CMD Vel Safety\\Node\\(50 Hz)};
        
        % Diff Drive Controller
        \node[draw, fill=red!20, rectangle, minimum width=3cm, minimum height=1.2cm] (diff_ctrl) at (0, -1) {Diff Drive\\Controller\\(50 Hz)};
        
        % Gazebo
        \node[draw, fill=purple!20, rectangle, minimum width=3cm, minimum height=1.5cm] (gazebo) at (0, -3) {Gazebo Sim\\(Physics)};
        
        % Beta EKF
        \node[draw, fill=orange!20, rectangle, minimum width=2.5cm, minimum height=1cm] (beta_ekf) at (5, 2) {Beta EKF\\(50 Hz)};
        
        % Dynamic Footprint
        \node[draw, fill=cyan!20, rectangle, minimum width=2.5cm, minimum height=1cm] (footprint) at (5, 0) {Dynamic\\Footprint\\(10 Hz)};
        
        % EKF Localization
        \node[draw, fill=pink!20, rectangle, minimum width=2.5cm, minimum height=1cm] (ekf_loc) at (5, -2) {EKF\\Localization\\(30 Hz)};
        
        % Topics
        \node[draw, ellipse, fill=gray!10] (cmd_nav2) at (-3, 4) {/cmd\_vel\_nav2};
        \node[draw, ellipse, fill=gray!10] (cmd_raw) at (-3, 2) {/cmd\_vel\_raw};
        \node[draw, ellipse, fill=gray!10] (cmd) at (-3, 0) {/cmd\_vel};
        \node[draw, ellipse, fill=gray!10] (beta) at (5, 4) {/trailer/beta};
        \node[draw, ellipse, fill=gray!10] (odom) at (5, -4) {/odom};
        
        % Arrows
        \draw[->, thick] (nav2) -- (cmd_nav2);
        \draw[->, thick] (cmd_nav2) -- (trailer_ctrl);
        \draw[->, thick] (trailer_ctrl) -- (cmd_raw);
        \draw[->, thick] (cmd_raw) -- (safety);
        \draw[->, thick] (safety) -- (cmd);
        \draw[->, thick] (cmd) -- (diff_ctrl);
        \draw[->, thick] (diff_ctrl) -- (gazebo);
        
        \draw[->, thick, dashed] (beta_ekf) -- (beta);
        \draw[->, thick, dashed] (beta) -- (trailer_ctrl);
        \draw[->, thick, dashed] (beta) -- (safety);
        \draw[->, thick, dashed] (beta) -- (footprint);
        
        \draw[->, thick, dashed] (gazebo) -- (odom);
        \draw[->, thick, dashed] (odom) -- (ekf_loc);
        \draw[->, thick, dashed] (odom) -- (beta_ekf);
        
        % Labels
        \node at (-1.5, 4) {\footnotesize Twist};
        \node at (-1.5, 2) {\footnotesize Twist};
        \node at (-1.5, 0) {\footnotesize Twist};
        \node at (6, 4) {\footnotesize Float64};
        \node at (6, -4) {\footnotesize Odometry};
    \end{tikzpicture}
    \caption{Architecture logicielle ROS2 du prototype}
    \label{fig:architecture_ros2}
\end{figure}

\subsection{Description des nœuds principaux}
\label{subsec:description_noeuds}

\begin{description}
    \item[Trailer-Aware Controller] Transforme les commandes Nav2 en commandes sécurisées pour la remorque en utilisant la cinématique inverse. Il réduit la vitesse angulaire lorsque l'angle d'attelage approche la limite et applique un terme correctif proportionnel pour aligner la remorque avec le tracteur.
    
    \item[CMD Vel Safety Node] Applique les limitations de vitesse finales pour prévenir le jackknife. Il calcule une vitesse angulaire maximale adaptative en fonction de l'angle $\beta$ courant et publie un statut de sécurité (0=safe, 1=warning, 2=danger).
    
    \item[Beta EKF Node] Estime l'angle d'attelage $\beta$ en fusionnant la mesure directe du joint hitch et une estimation cinématique dérivée des vitesses des roues. Il publie également le rayon de virage minimal et la position de l'ICR (Instantaneous Center of Rotation).
    
    \item[Dynamic Footprint Node] Recalcule en temps réel l'empreinte de collision articulée du système tracteur-remorque sous forme de multi-cercles, permettant aux costmaps Nav2 de prendre en compte la géométrie variable de la remorque.
    
    \item[EKF Localization] Gère la double EKF (locale et globale) pour la fusion odométrie/IMU/GPS, publiant les transformations \texttt{map$\to$odom} et \texttt{odom$\to$base\_footprint}.
\end{description}

\section{Algorithmes implémentés pour la navigation autonome}
\label{sec:algorithmes_navigation}

\subsection{Fusion de capteurs pour estimation de $\beta$}
\label{subsec:fusion_beta}

L'estimation de l'angle d'attelage utilise un EKF fusionnant deux sources :

\begin{table}[htbp]
    \centering
    \rowcolors{2}{gray!10}{white}
    \begin{tabular}{p{3cm}p{4cm}p{4cm}}
        \toprule
        \textbf{Étape} & \textbf{Opération} & \textbf{Équation} \\
        \midrule
        Prédiction & Modèle cinématique & $\dot{\beta} = \omega - \frac{v}{d}\sin\beta$ \\
        Covariance & Jacobien transition & $P = F P F^T + Q \Delta t$ \\
        Mise à jour joint & Mesure hitch & $K = P H^T (H P H^T + R_{\text{joint}})^{-1}$ \\
        Mise à jour roues & Estimation cinématique & $\beta_{\text{wheel}} = \arctan(\omega d / v)$ \\
        Limitation & Bornes physiques & $\beta \in [-\beta_{\max}, \beta_{\max}]$ \\
        \bottomrule
    \end{tabular}
    \caption{Étapes de l'EKF pour estimation de $\beta$}
    \label{tab:ekf_beta}
\end{table}

\subsection{Détection et prévention du jackknife}
\label{subsec:anti_jackknife}

La prévention du jackknife repose sur une détection de zone et une limitation adaptative :

\begin{table}[htbp]
    \centering
    \rowcolors{2}{gray!10}{white}
    \begin{tabular}{lccc}
        \toprule
        \textbf{Zone} & \textbf{Condition} & \textbf{Réduction $\omega_{\max}$} & \textbf{Statut} \\
        \midrule
        Safe & $|\beta| < 30^\circ$ & $\omega_{\max,adaptive}$ & 0 \\
        Warning & $30^\circ \leq |\beta| < 40^\circ$ & $0.7 \times \omega_{\max,adaptive}$ & 1 \\
        Danger & $|\beta| \geq 40^\circ$ & $0.5 \times \omega_{\max,adaptive}$ & 2 \\
        \bottomrule
    \end{tabular}
    \caption{Zones de sécurité et limitations associées}
    \label{tab:zones_jackknife}
\end{table}

\subsection{Contrôle du recul sécurisé}
\label{subsec:controle_recul}

Le contrôle du recul combine correction inverse de cinématique et limitation adaptative :

\begin{table}[htbp]
    \centering
    \rowcolors{2}{gray!10}{white}
    \begin{tabular}{p{2.5cm}p{5cm}p{4cm}}
        \toprule
        \textbf{Étape} & \textbf{Opération} & \textbf{Formule} \\
        \midrule
        1. Inverse cinématique & Réduction courbure si $\beta_{\text{steady}} > \beta_{\max}$ & $\kappa_{\max} = \tan(\beta_{\max})/d$ \\
        2. Correction alignement & Terme proportionnel si $|\beta| > 0.1$ rad & $\omega_{\text{corr}} = 0.5(v/d)\sin\beta$ \\
        3. Réduction vitesse & Scale si $|\beta| > 0.2$ rad & $\text{scale} = 1-0.5|\beta|/\beta_{\max}$ \\
        4. Limite finale & Bornes cinématiques & $\omega_{\max} = \frac{\sin\beta_{\max}}{d}|v|$ \\
        \bottomrule
    \end{tabular}
    \caption{Étapes du contrôle du recul sécurisé}
    \label{tab:controle_recul}
\end{table}

\section{Description du prototype Gazebo}
\label{sec:prototype_gazebo}

\subsection{Choix des primitives géométriques}
\label{subsec:choix_primitives}

Le prototype utilise des primitives géométriques simples (boîtes) pour validation rapide :

\begin{table}[htbp]
    \centering
    \rowcolors{2}{gray!10}{white}
    \begin{tabular}{lp{10cm}}
        \toprule
        \textbf{Avantage} & \textbf{Description} \\
        \midrule
        Rapidité & Définition URDF simple, pas de CAO détaillée \\
        Validation théorique & Tests équations cinématiques et algorithmes fusion \\
        Performance & Calcul physique réduit, tests rapides \\
        Flexibilité & Ajustement paramètres géométriques facile \\
        \bottomrule
    \end{tabular}
    \caption{Avantages des primitives géométriques pour le prototypage}
    \label{tab:avantages_primitives}
\end{table}

Le modèle URDF définit le tracteur comme un châssis principal avec quatre roues et la remorque comme un corps articulé relié par un joint d'attelage de type \texttt{revolute} avec des limites $\pm 45^\circ$.

\subsection{Configuration des capteurs simulés}
\label{subsec:config_capteurs}

\begin{table}[htbp]
    \centering
    \rowcolors{2}{gray!10}{white}
    \begin{tabular}{lcccc}
        \toprule
        \textbf{Capteur} & \textbf{Portée/Specs} & \textbf{Fréquence} & \textbf{Topic} \\
        \midrule
        LIDAR & 0.15-12.0 m, $\pm 120^\circ$ & 10 Hz & /scan \\
        IMU & Gyro+Accéléro bruit gaussien & 50 Hz & /imu \\
        GPS & $\sigma_{\text{pos}} = 0.5$ m & 10 Hz & /gps/fix \\
        Caméra & 640$\times$480, $90^\circ$ FOV & 30 Hz & /image\_raw \\
        \bottomrule
    \end{tabular}
    \caption{Configuration des capteurs simulés dans Gazebo}
    \label{tab:config_capteurs}
\end{table}

\subsection{Environnement de simulation}
\label{subsec:environnement_sim}

Le monde \texttt{warehouse.world} définit un environnement d'entrepôt avec obstacles statiques (murs, étagères, caisses) pour tester la navigation et l'évitement d'obstacles. Position spawn par défaut : $(x=-0.14, y=-2.07)$ m.

\section*{Conclusion}
\addcontentsline{toc}{section}{Conclusion}

Ce chapitre a présenté l'étude théorique complète du système tracteur-remorque :

\begin{table}[htbp]
    \centering
    \rowcolors{2}{gray!10}{white}
    \begin{tabular}{lp{10cm}}
        \toprule
        \textbf{Apport} & \textbf{Description} \\
        \midrule
        Équation cinématique & $\dot{\beta} = \omega - \frac{v}{d}\sin\beta$ : fondement théorique \\
        Rayon minimal & $R_{\min} \approx 0.71$ m : limite de stabilité jackknife \\
        Architecture EKF & Fusion robuste odométrie/IMU/GPS à 3 niveaux \\
        Contrôle sécurisé & Prévention jackknife par limitations adaptatives \\
        Validation prototype & Gazebo primitives géométriques pour tests rapides \\
        \bottomrule
    \end{tabular}
    \caption{Synthèse des apports du chapitre}
    \label{tab:synthese_chapitre}
\end{table}

\textbf{Limites actuelles :} absence de modélisation détaillée des frottements d'attelage, dynamique remorque simplifiée (modèle cinématique pur).

\textbf{Transition Chapitre 2 :} passage à la modélisation CAO détaillée et intégration sur plateforme physique Husky.

\end{document}
